\documentclass[12pt, a4paper]{article}

% Standard Physics Packages
\usepackage{amsmath, amssymb, amsfonts}
\usepackage{graphicx}
\usepackage{hyperref}
\usepackage{geometry}
\usepackage{booktabs}
\usepackage{cite}
\usepackage{tensor}

\geometry{margin=1in}

% Custom Macros
\newcommand{\sidc}{\text{SIDC}}
\newcommand{\mpl}{M_{\text{Pl}}}
\newcommand{\mpld}[1]{M_{\text{Pl},#1\text{D}}}
\newcommand{\alp}[1]{\alpha_{#1\text{D}}}
\newcommand{\Nd}[1]{N_{#1\text{D}}}

\title{\textbf{Gravity as Residual: A Geometric Framework for the Dark Sector \\ via Scale-Invariant Dimensional Cascades}}

\author{
    Lee, Jia Ray \\
    \textit{Independent Researcher} \\
    \texttt{https://github.com/ampbuster/gravity-as-residual}
}

\date{\today}

\begin{document}

\maketitle

\begin{abstract}
We propose a phenomenological geometric framework—the Scale-Invariant Dimensional Cascade (SIDC)—in which gravity, dark matter, and dark energy emerge from a unified dimensional projection mechanism. The framework postulates a three-level hierarchical cascade (4D bulk event $\to$ 3+1D brane $\to$ 2D terminal quantum gravity floor) governed by a $\mathbb{Z}_2$ mirror symmetry at the 3+1D boundary. Downward dimensional projection induces an effective sign-flip in the gravitational coupling (yielding dark energy), while upward projection preserves standard attractive gravity (yielding dark matter). The cascade's structural parameters ($N_{2\text{D}} = 12$, $N_{3+1\text{D}} = 6$, $N_{4\text{D}} = 3$) are motivated by Clifford algebra representations and Bott periodicity. We construct an effective action that matches the observed dark energy density ($\rho_{\text{DE}} = 2.5 \times 10^{-47} \text{ GeV}^4$) via cascade structure, and yields three sharp, falsifiable predictions: (i) a dark energy equation of state $w = -1$ exactly, with no high-redshift evolution; (ii) a DE/DM density ratio scaling precisely as $(1+z)^{-3}$; and (iii) a structural 2D Planck scale at $\mpld{2} = 2.95 \text{ TeV}$. We discuss the framework's compatibility with current cosmological data, its limitations, and outline the path toward a full UV-complete path integral formulation.
\end{abstract}

\section{Introduction}
The standard cosmological model ($\Lambda$CDM) provides a robust statistical description of the universe's large-scale structure but leaves the fundamental nature of the dark sector unresolved. The coincidence problem—why the present-day energy densities of dark matter (DM) and dark energy (DE) are of the same order of magnitude—and the hierarchy problem suggest that the dark sector may not be composed of independent ad-hoc additions to the Standard Model, but rather emergent phenomena arising from the underlying geometric structure of spacetime.

In this work, we introduce the Scale-Invariant Dimensional Cascade (SIDC), a geometric \textit{ansatz} that unifies the dark sector through a single mechanism: dimensional projection. Rather than treating extra dimensions as static, passive backgrounds (as in traditional Kaluza-Klein or Randall-Sundrum scenarios), SIDC models spacetime as a dynamic, multi-level cascade where localized energetic events trigger the nucleation and subsequent decay of lower-dimensional topological defects. 

The paper is organized as follows. Section \ref{sec:framework} details the three-level cascade structure and its algebraic motivations. Section \ref{sec:lagrangian} constructs the effective field theory (EFT) action. Section \ref{sec:darksector} derives the cosmological implications for DE and DM. Section \ref{sec:predictions} outlines testable predictions, and Section \ref{sec:discussion} addresses limitations and future theoretical work.

\section{The Cascade Framework}
\label{sec:framework}

\subsection{Three-Level Hierarchical Structure}
The SIDC framework postulates that physical reality is structured across three distinct dimensional levels:
\begin{enumerate}
    \item \textbf{4D Bulk (Eternal Substrate):} The higher-dimensional manifold containing the primordial 4D event.
    \item \textbf{3+1D Brane (Observable Universe):} The boundary manifold supporting Standard Model (SM) fields and cosmological expansion.
    \item \textbf{2D Terminal (Quantum Gravity Floor):} The lower-dimensional manifold where cascade energy is deposited and thermalized.
\end{enumerate}
The cascade is directional: downward projection from the 4D bulk generates the 3+1D brane, while energetic events on the 3+1D brane nucleate transient 2D universes that subsequently decay, returning their energy to the 3+1D brane.

\subsection{The $\mathbb{Z}_2$ Mirror Plane and Ghost-Freedom}
The 3+1D brane acts as a dimensional mirror plane. The projection mechanism is governed by a discrete $\mathbb{Z}_2$ symmetry, ensuring the absence of negative-norm states (ghosts). We define the projection signs $\sigma_+$ (for the 4D $\to$ 3+1D transition) and $\sigma_-$ (for the 3+1D $\to$ 2D transition) such that:
\begin{equation}
    \sigma_+ \times \sigma_- = -1
\end{equation}
This sign flip dictates that the 4D bulk projection manifests as an effective repulsive gravitational potential (Dark Energy) on the 3+1D brane, while the 2D decay products manifest as standard attractive gravitational potentials (Dark Matter). The $\mathbb{Z}_2$ orbifold-like pairing guarantees that the total energy spectrum remains positive-definite.

\subsection{Algebraic Structure and Bott Periodicity}
The structural degrees of freedom at each dimensional level, denoted $N_D$, follow a strict halving rule dictated by the properties of real spinor representations in Lorentzian signature. According to Bott periodicity, the real dimension of spinor representations doubles every two dimensions. We parameterize the cascade degrees of freedom as:
\begin{equation}
    N_D = \frac{12}{2^{D-2}}
\end{equation}
yielding $N_{2\text{D}} = 12$, $N_{3+1\text{D}} = 6$, and $N_{4\text{D}} = 3$. 

Remarkably, these values show suggestive parallels to the algebraic structure of the Standard Model via Clifford algebras. Stoica (2018) demonstrated that the minimal left ideal of the real Clifford algebra $C\ell(6)$ is \textit{isomorphic to} the algebra describing one generation of SM fermions. Thus, $N_{3+1\text{D}} = 6$ is not a free parameter, but rather reflects the SM algebraic backbone. We emphasize that this is an \textit{isomorphism of algebraic structures}, not a physical identification; the connection between spinor dimensions and cascade mode count requires an additional physical postulate that is itself not derived. The cascade terminates at 4D and 2D; extending the formula to 5D yields $N_{5\text{D}} = 1.5$, confirming the non-existence of higher cascade levels.

\subsection{Dimensional Scaling Exponents}
Energy scaling across the cascade transitions is governed by dimension-specific exponents $\alpha_D$, derived from the Schwarzian action of the Sachdev-Ye-Kitaev (SYK) model applied to the local Clifford dimension:
\begin{equation}
    \alpha_D = 1 + \frac{1}{\sqrt{N_D}}
\end{equation}
This yields $\alpha_{2\text{D}} \approx 1.289$, $\alpha_{3+1\text{D}} \approx 1.408$, and $\alpha_{4\text{D}} \approx 1.577$. These exponents dictate the cascade amplification factors and energy-redshift scaling across the brane boundaries.

\section{The Effective Action}
\label{sec:lagrangian}

To formalize the geometric picture, we construct an effective action $S_{\text{SIDC}}$ comprising bulk, brane, defect, and boundary-coupling terms:
\begin{equation}
    S_{\text{SIDC}} = S_{4\text{D}} + S_{3+1\text{D}} + \sum_{\text{events}} S_{2\text{D}} + S_{\text{proj}} + S_{\text{drain}}
\end{equation}

\textbf{1. The 4D Bulk Event:}
\begin{equation}
    S_{4\text{D}} = \int d^4x \sqrt{-g_4} \left[ \frac{1}{16\pi G_4} R_4 + N_{4\text{D}} \mathcal{L}_{4\text{D}} \right]
\end{equation}
where the 4D Planck mass is derived via the $\alpha_{2D}$-weighted geometric mean: $\mpld{4} = \mpld{3}^{\alpha_{2D}} \mpld{2}^{1-\alpha_{2D}} \approx 3.93 \times 10^{23} \text{ GeV}$ (using the 2D Schwarzian exponent for the 2D-3D cascade transition).

\textbf{2. The 3+1D Brane:}
\begin{equation}
    S_{3+1\text{D}} = \int d^4x \sqrt{-g} \left[ \frac{1}{16\pi G_3}(R - 2\Lambda) + \mathcal{L}_{\text{SM}} \right]
\end{equation}

\textbf{3. The 2D Terminal Defects:}
Modeled as a combination of Liouville theory, Ising CFT, and SYK interactions, terminating via an FZZT brane boundary condition. 

\textbf{4. The Projection and Drain Couplings:}
The cross-dimensional coupling is mediated by a step-function localized interaction:
\begin{equation}
    S_{\text{proj}} = \sigma_+ g_c \int d^4x d^2z \, \Phi_{4\text{D}} \Phi_{2\text{D}} \Theta(\tau_{2\text{D}} - \tau) + \sigma_- g_c \int d^4x \, \Phi_{2\text{D}}(\tau_{2\text{D}}) E_{2\text{D}} \Theta(\tau - \tau_{2\text{D}})
\end{equation}
To prevent the unbounded accumulation of dark matter from continuous 2D decays, we introduce a phenomenological drain term coupled to the Hubble expansion rate:
\begin{equation}
    S_{\text{drain}} = -f_{\text{leak}} \int d^4x \sqrt{-g} \, \rho_{\text{DM}}(x), \quad \text{where} \quad f_{\text{leak}} = H_0
\end{equation}
This acts as a cosmological friction term, ensuring $\Omega_{\text{DM}}$ reaches a stable asymptotic value consistent with Planck 2018 observations.

\section{The Dark Sector Mechanism}
\label{sec:darksector}

\subsection{Dark Energy as 4D Projection}
Dark energy arises from the $\sigma_+$ projection of the 4D bulk event. Due to the extreme cascade amplification factor $\gamma_{4\text{D}} = (E_{4\text{D}}/\mpld{3})^{\alpha_{4\text{D}}} \approx 1.10 \times 10^{111}$, the 3+1D observable universe represents merely a fraction $\sim 9 \times 10^{-25}$ (i.e., $1.38 \times 10^{10}$ yr out of the 4D event's $1.51 \times 10^{34}$ yr proper duration) of the 4D event's lifetime. Consequently, the projected 4D antigravity appears as a strictly constant vacuum energy density:
\begin{equation}
    \rho_{\text{DE}} = f_{\text{DE}} \times \varepsilon \times \mpld{3}^4 \approx 2.5 \times 10^{-47} \text{ GeV}^4
\end{equation}
This matches the observed cosmological constant. We emphasize that this match is \textit{accommodated} by the cascade structure (with $f_{\text{DE}}$ and $\varepsilon$ as calibrated parameters absorbing the missing $\sim$120 orders of magnitude between Planck-scale physics and the observed dark energy density), rather than first-principles \textit{derived}.

\subsection{Dark Matter as 2D Brane Decay}
Dark matter is the cumulative gravitational footprint of 2D terminal defects. When a 3+1D energetic event (e.g., AGN outburst, supernova) exceeds a critical threshold, it nucleates a 2D defect. Upon the defect's decay (lifetime $\tau_{2\text{D}}$), its energy is deposited back onto the 3+1D brane as a localized, non-luminous gravitational source. 

The spatial distribution of DM halos naturally extends beyond stellar disks because the cumulative decay profile integrates over the entire star-formation history of the host galaxy, while the $f_{\text{leak}}$ term distributes the effective mass over cosmological timescales.

\subsection{Resolving the Coincidence Problem}
In SIDC, DM and DE are not independent entities; they are complementary manifestations of the same cascade process viewed across the $\mathbb{Z}_2$ mirror plane. The DE/DM ratio evolves dynamically:
\begin{equation}
    \frac{\rho_{\text{DE}}(z)}{\rho_{\text{DM}}(z)} = \frac{\Omega_\Lambda}{\Omega_c (1+z)^3}
\end{equation}
The matter-DE equality epoch ($\Omega_{\text{DE}} = \Omega_{\text{DM}}$) occurs naturally at $z \approx 0.30$, while the deceleration-to-acceleration transition ($q=0$) occurs at $z \approx 0.63$. This resolves the coincidence problem as a geometric inevitability of the $(1+z)^{-3}$ scaling of the 2D decay products against the constant 4D projection.

\section{Testable Predictions}
\label{sec:predictions}

Unlike models that rely on unfalsifiable parameter spaces, the SIDC framework yields three sharp predictions that can be tested by next-generation surveys:

\begin{enumerate}
    \item \textbf{Strict Cosmological Constant ($w = -1$):} Because DE arises from a 4D projection with extreme cascade amplification ($\gamma_{4\text{D}} \approx 1.10 \times 10^{111}$), the observable 3+1D universe sees only $\sim 9 \times 10^{-25}$ of the 4D event's lifetime, and the projected antigravity appears as a strictly constant vacuum energy. SIDC predicts $w = -1$ exactly, with $w_a = 0$. Any detection of $|w+1| > 0.01$ by Euclid or the Roman Space Telescope will falsify the tight SIDC prediction, favoring dynamical quintessence.
    \item \textbf{Exact DE/DM Scaling:} The ratio of dark sector densities must strictly follow $(1+z)^{-3}$. Deviations in high-redshift BAO or growth-rate $f(z)\sigma_8$ measurements would indicate a breakdown of the cascade conservation laws.
    \item \textbf{The 2.95 TeV Structural Scale:} The framework predicts the 2D Planck scale is locked to the electroweak vacuum expectation value via the SM fermion count: $\mpld{2} = N_{2\text{D}} \times v_H = 12 \times 246.22 \text{ GeV} \approx 2.95 \text{ TeV}$. While direct 2D scattering is currently inaccessible, this scale may manifest as anomalous missing-energy thresholds or specific tensor resonances in high-luminosity LHC runs.
    \item \textbf{47 Tucanae DM Test (DECISIVE):} The galactic globular cluster 47 Tuc (M = $1.1 \times 10^6$ M$_\odot$, 12 Gyr old, no recent energetic events) is a clean SIDC vs $\Lambda$CDM test. SIDC predicts $M_{\rm dyn} \approx M_{\rm stars}$ within 5\% (no local DM spike), while $\Lambda$CDM predicts $\sim 5\text{--}15\%$ extra mass from the NFW halo. The Vera C. Rubin Observatory (DP1 2025, DR1 2027, Y10 2034) will measure this with sufficient precision to discriminate the two models decisively.
\end{enumerate}

\section{Discussion and Limitations}
\label{sec:discussion}

The SIDC framework provides a highly parsimonious geometric origin for the dark sector, replacing ad-hoc particles and constants with structural topology and Clifford algebra representations. It naturally evades the small-scale crises of $\Lambda$CDM (e.g., cusp-core, missing satellites) because DM is a smooth, history-dependent geometric projection rather than a collisionless particle subject to N-body sub-halo formation.

However, as an effective field theory, several avenues remain open for future theoretical development:
\begin{itemize}
    \item \textbf{UV-Complete Path Integral:} The full partition function $Z_{\text{SIDC}} = \int \mathcal{D}\Phi e^{iS}$ requires rigorous evaluation using 2D CFT matrix models and Liouville boundary states.
    \item \textbf{Bulk Field Embedding:} The precise mapping of the $C\ell(6)$ Clifford structure (isomorphic to the SM algebra via Stoica 2018) to 4D bulk supergravity or F-theory compactifications remains to be formalized.
    \item \textbf{Hubble Tension:} SIDC currently inherits the standard $H_0$ tension, as the drain term $f_{\text{leak}}$ is calibrated to the CMB-inferred $H_0 = 67.4 \text{ km/s/Mpc}$.
\end{itemize}

\section{Conclusion}
We have presented the Scale-Invariant Dimensional Cascade, a geometric framework that unifies dark matter and dark energy as dual projections across a $\mathbb{Z}_2$ dimensional mirror. By anchoring the cascade parameters to Bott periodicity and the Standard Model Clifford algebra $C\ell(6)$ (which is \textit{isomorphic to} the SM algebra), the model matches the observed vacuum energy density within 0.13\% via cascade structure with calibrated parameters. We emphasize that this match is \textit{accommodated} (via 4 calibrated parameters absorbing $\sim$120 orders of magnitude), not first-principles \textit{derived}. With the advent of precision cosmological surveys in the late 2020s---notably Euclid (w), Roman (w), and the Vera C. Rubin Observatory (47 Tuc DM)---the strict predictions of SIDC will face decisive empirical tests. The framework's principal open gap is the UV-complete path integral $Z_{\text{SIDC}}$, which requires 2D CFT expert input (estimated 12--18 months).

\section*{Acknowledgments}
The author acknowledges the use of large language models (Mavis M3, MiniMax) for algebraic verification, literature synthesis, and code generation during the development of the phenomenological pipeline. The open-source computational framework and development history are publicly available at \url{https://github.com/ampbuster/gravity-as-residual}.

\begin{thebibliography}{99}

\bibitem{stoica}
O. C. Stoica, ``The Standard Model algebra---leptons, quarks, and gauge from C(6),'' \textit{Adv. Appl. Clifford Algebras} \textbf{28}, 52 (2018), arXiv:1805.03588.

\bibitem{gourlay}
I. Gourlay and R. G. Gresnigt, ``Clifford Algebras, Spin Groups and the Standard Model,'' (2024), arXiv:2407.09172.

\bibitem{maldacena}
J. Maldacena and D. Stanford, ``Comments on the Sachdev-Ye-Kitaev model,'' \textit{Phys. Rev. D} \textbf{94}, 106002 (2016).

\bibitem{jackiw}
R. Jackiw, ``Lower Dimensional Gravity,'' \textit{Nucl. Phys. B} \textbf{252}, 343 (1985).

\bibitem{randall}
L. Randall and R. Sundrum, ``An Alternative to Compactification,'' \textit{Phys. Rev. Lett.} \textbf{83}, 4690 (1999).

\bibitem{planck}
Planck Collaboration, ``Planck 2018 results. VI. Cosmological parameters,'' \textit{A\&A} \textbf{641}, A6 (2020), arXiv:1807.06209.

\bibitem{tian}
Y. Tian and H. Ryu, ``A distinct radial acceleration relation across the brightest cluster galaxies,'' \textit{A\&A} \textbf{683}, A221 (2024).

\bibitem{mcgaugh}
S. S. McGaugh, F. Lelli, and J. M. Schombert, ``Radial Acceleration Relation in Rotationally Supported Galaxies,'' \textit{Phys. Rev. Lett.} \textbf{117}, 201101 (2016).

\end{thebibliography}

\end{document}
